% !TEX program = pdflatex
\documentclass[11pt,a4paper]{article}
\usepackage{amsmath,amssymb,amsthm,amsfonts}
\usepackage{geometry}
\geometry{margin=1in}

% Theorem environments
\theoremstyle{definition}
\newtheorem{definition}{Definition}[section]
\newtheorem{theorem}{Theorem}[section]
\newtheorem{lemma}{Lemma}[section]
\newtheorem{corollary}{Corollary}[section]
\newtheorem{proof_outline}{Proof Sketch}

\title{A Formulation of the Truth-Preservation Incompleteness Theorem}
\author{Masaomi Chiba}
\date{2026-04-29}

\begin{document}

\maketitle

\begin{abstract}
This paper generalizes the concept of ``truth preservation'' in formal and dynamical systems and formulates a metatheorem showing that, for any truth preservation system $T$, there necessarily exists a proposition $G_T$ that is unprovable within $T$. The theorem subsumes G\"{o}del's incompleteness theorems and Chaitin's incompleteness theorem in a unified manner, interpreting them as structural limitations arising from the \emph{a priori} requirement of self-identity and invariant preservation.
\end{abstract}

\section{Introduction}

Mathematical, logical, and physical systems are constructed on the premise that certain states, relations, or invariants are maintained and preserved throughout processes of inference or temporal evolution. In this paper, we generalize this property and refer to it as \emph{truth preservation}. The aim of this paper is to formally describe how, whenever a system requires some form of truth preservation, structural unprovability inevitably arises from the very foundation of that preservation law.

\section{Definitions}

\begin{definition}[Truth Preservation System $T$]
A \emph{Truth Preservation System} $T$ is a pair consisting of a state space $X$ and a collection of operations on $X$ (inference rules, time-evolution operators, etc.), defined so as to permanently maintain, through any finite sequence of operations, a specific property $P$ assigned to the initial state---such as an invariant, consistency, or self-identity $A = A$.
\end{definition}

\begin{definition}[Provability in $T$]
A proposition $S$ is said to be \emph{provable under $T$} if $S$ is derivable by a finite sequence of deductive inferences or computational procedures that presuppose the truth preservation property $P$ required by $T$. This is written $T \vdash S$.
\end{definition}

\section{Main Theorem}

\begin{theorem}[Truth-Preservation Incompleteness Theorem]
\label{thm:incompleteness}
For any sufficiently powerful truth preservation system $T$, there necessarily exists a proposition $G_T$ that is expressible in the language of $T$ but can be neither proved nor disproved under $T$:
\[
\forall T,\; \exists G_T \;\text{ s.t. }\; (T \nvdash G_T) \land (T \nvdash \neg G_T).
\]
\end{theorem}

\begin{proof_outline}
The structural reason for Theorem~\ref{thm:incompleteness} is as follows.
\begin{enumerate}
    \item In order for $T$ to carry out any inference or computation, the truth preservation property $P$ must be given \emph{a priori} as a precondition.
    \item If system $T$ attempts to prove its own completeness, it must derive---using $T$'s internal operations---either ``$P$ is always maintained'' or ``the boundary conditions that specify $P$ are valid.''
    \item However, the internal operations of $T$ already depend, by definition, on the validity of $P$. Consequently, when $T$ attempts to evaluate any proposition $G_T$ concerning the limits of $P$'s validity or the global topological properties generated by iterated truth-preserving operations (e.g., the limits of unbounded extension beyond the system), $T$ falls into circular reasoning or becomes undecidable.
    \item Therefore, as long as $T$ is held fixed, it is structurally impossible to determine within $T$ the boundary property $G_T$ that $P$ itself gives rise to.
\end{enumerate}
\end{proof_outline}

\section{Subsumption of Known Incompleteness Theorems}

The present theorem allows the historically known incompleteness theorems to be positioned as concrete instances of a truth preservation system $T$.

\begin{corollary}[Subsumption of G\"{o}del's First Incompleteness Theorem]
Let $T_{\text{G\"{o}del}}$ be a formal arithmetic system whose truth preservation property is \emph{logical consistency} (the non-derivability of falsehood). Then the self-referential proposition $G_{\text{G\"{o}del}}$, which asks about the limits of the logical truth preservation of $T_{\text{G\"{o}del}}$, is unprovable within $T_{\text{G\"{o}del}}$.
\end{corollary}

\begin{corollary}[Subsumption of Chaitin's Incompleteness Theorem]
Let $T_{\text{Chaitin}}$ be a system whose invariant of truth preservation is the \emph{algorithmic information content} (Kolmogorov complexity) of the axiom system. Then the proposition $G_{\text{Chaitin}}$, which concerns randomness exceeding the information content of $T_{\text{Chaitin}}$, cannot be compressed (proved) by the truth-preserving operations internal to $T_{\text{Chaitin}}$.
\end{corollary}

\section{Conclusion}

The act of ``proof'' is a local operation that depends on a truth preservation system $T$. The requirement of truth preservation---self-identity and logical consistency---is at once the indispensable driving condition of a system and the inevitable source of a structural blind spot ($G_T$) that is unobservable from within that system. The present theorem demonstrates that incompleteness holds universally, beyond the concept of self-identity, in every system that requires any invariant whatsoever.

\end{document}
